\subsection{Families of Features}
\label{sec:feature-types}
\slidesrc{Text features: 2--13}

\subsubsection{Stylistic features}
\slidesrc{Text features: 2--8}

\textbf{Stylistic} features describe \emph{how} something is written, not \emph{what}
it is about: sentence length, punctuation, the proportion of pronouns or adjectives.
Two people writing about the same topic use different styles, and the same person
writing about different topics keeps much of their style. This is why stylistic
features are the basis of \textbf{author profiling} (guessing the age, gender, native
language, personality of the author) and \textbf{authorship attribution} (who wrote this
text?).

The slides show the 29 stylistic features used by an author profiling system for four
languages:

\begin{center}
\footnotesize
\begin{tabularx}{\linewidth}{@{}r X c | r X c@{}}
\toprule
\textbf{No.} & \textbf{feature} & \textbf{all 4} & \textbf{No.} & \textbf{feature} &
  \textbf{all 4} \\
\midrule
1 & \% question sentences & \checkmark & 16 & \% words with 3+ syllables & \\
2 & \% short sentences & \checkmark & 17 & average syllables per word & \\
3 & \% long sentences & \checkmark & 18 & \% adjectives & \\
4 & average sentence length & \checkmark & 19 & \% adverbs & \\
5 & average word length & \checkmark & 20 & \% capitals & \checkmark \\
6 & \% words with 6+ letters & \checkmark & 21 & \% colons & \checkmark \\
7 & \% words with 2--3 letters & \checkmark & 22 & \% determiners & \\
8 & \% semicolons & \checkmark & 23 & \% digits & \checkmark \\
9 & \% punctuation & \checkmark & 24 & \% full stops & \checkmark \\
10 & \% pronouns & & 25 & \% interjections & \\
11 & \% prepositions & & 26 & \% modals & \\
12 & \% coordinating conjunctions & & 27 & \% nouns & \\
13 & \% commas & \checkmark & 28 & \% personal pronouns & \\
14 & \% articles & & 29 & \% verbs & \\
15 & \% words with one syllable & & & & \\
\bottomrule
\end{tabularx}
\captionof{table}{The stylistic features of an author profiling system (from a paper
cited in the slides). All 29 are used for English; ``all 4'' marks those used also for
Dutch, Spanish and Italian. The others need a POS tagger or a syllable counter, which
the system had only for English.}
\label{tab:stylistic-features}
\end{center}

\begin{nbintuition}[why only some features in every language?]
Features 1--9, 13, 20, 21, 23, 24 can be computed by simply counting characters, words
and sentences: they work in any language. Features like ``\% adjectives'' need a POS
tagger (Section~\ref{sec:pos-tagging}) and ``syllables per word'' needs a syllable
counter, which are language-specific tools. A recurring lesson: the cheapest features
are the most portable.
\end{nbintuition}

\paragraph{N-grams and collocations.}
The co-occurrence of the same $n$ characters or words: character n-grams (\emph{ing\_},
\emph{\_th}) capture spelling and morphology habits, word n-grams and collocations
(\emph{strong tea}, Section~\ref{sec:corpora}) capture habitual phrasing.

\paragraph{POS (part of speech).}
Playing a particular role (``part'') in a sentence (``speech'') is a feature of
\emph{one word inside a sentence}: no need for the full document or corpus. The
\emph{number} of words with a given POS (number of pronouns, of adjectives, \ldots),
possibly normalised by the length of the document, is a feature of the
\emph{document}.

\paragraph{FW (function words).}

\begin{nbdef}{Function word}
In linguistics, \textbf{function words} (also called \emph{functors}) are words that
have little lexical meaning, or an ambiguous meaning, and express grammatical
relationships among other words within a sentence, or specify the attitude or mood of
the speaker.
\end{nbdef}

Being a function word is a feature of the \emph{word itself}: if the word is not
ambiguous, we do not even need the sentence it belongs to. The \emph{number} of function
words, possibly normalised by the document length, is a feature of the document. Kinds of
function words:
\begin{itemize}[nosep]
  \item \textbf{articles}: \emph{the}, \emph{a} (in some inflected languages the article
        takes the case of the following noun);
  \item \textbf{pronouns}: inflected in English (\emph{he}--\emph{him},
        \emph{she}--\emph{her});
  \item \textbf{conjunctions}: \emph{and} (uninflected in English);
  \item \textbf{subordinating conjunctions}: \emph{if}, \emph{then}, \emph{well},
        \emph{however}, \emph{thus};
  \item \textbf{auxiliary verbs}: part of the conjugation of main verbs, always
        inflected (\emph{has} gone, \emph{will} go);
  \item \textbf{particles}: \emph{up}, \emph{on}, \emph{down} (give \emph{up});
  \item \textbf{interjections}: sometimes called ``filled pauses'' (\emph{uh}, \emph{um});
  \item \textbf{expletives}: words that take the place of a sentence, among other
        functions;
  \item \textbf{pro-sentences}: \emph{yes}, \emph{okay}.
\end{itemize}
There is a \textbf{huge overlap with stop words} (Section~\ref{sec:stop-words-tfidf}):
the word cloud on the slide, of a novel, is dominated by \emph{the}, \emph{and},
\emph{to}, \emph{of}, \emph{a}, \emph{was}, \emph{her}, \emph{he}.

\begin{nbwarning}[stop words are not always noise]
For topic classification we \emph{remove} stop words because they say nothing about the
topic. For author profiling we \emph{count} them, because \emph{how often} an author uses
\emph{the}, \emph{of}, \emph{I}, \emph{she} is one of the best clues to who the author is.
The same words are noise for one task and signal for another.
\end{nbwarning}

\paragraph{Readability.}
\textbf{Readability} measures how easy a piece of text is to read. It can include
elements of complexity, familiarity, legibility and typography. Readability formulas
usually look at sentence length, syllable density and word familiarity; an example is
the Flesch--Kincaid grade level of Section~\ref{sec:course-map}. Online tools (e.g.\
\url{https://readable.com}) compute many such scores.

\subsubsection{Content-based features}
\slidesrc{Text features: 9}

\textbf{Content-based} features describe \emph{what} the text is about:
\begin{itemize}[nosep]
  \item the frequency of words belonging to some given, known \textbf{semantic category}
        (dictionary-based, topic-based): e.g.\ how many words are about food, sport,
        money;
  \item \textbf{emotion words} (based on Ekman's or Plutchik's basic emotions,
        Section~\ref{sec:course-map});
  \item \textbf{TF-IDF} (Section~\ref{sec:stop-words-tfidf}). Recall:
        \begin{align*}
          \text{tf-idf}(t, d, C) &= \text{tf}(t, d)\cdot\text{idf}(t, C),\\
          \text{tf}(t, d) &= \frac{\#\,t \text{ in } d}{\text{length}(d)}, \qquad
          \text{idf}(t, C) = \log\frac{\abs{C}}{\abs{\{d \in C : t \in d\}}};
        \end{align*}
  \item \textbf{LSA} (Latent Semantic Analysis, Section~\ref{sec:dsm-in-practice}):
        dense vectors of the words or of the document.
\end{itemize}

\begin{nbwarning}[LSA and word2vec]
The slide writes ``LSA, implemented by word2vec''. Strictly, they are two different
methods that produce the same kind of object, a dense vector per word: LSA reduces a
word--document matrix with SVD (Section~\ref{sec:dsm-in-practice}), word2vec trains a
small neural network to predict words from their contexts. Both are distributional, and
both can be used as content features.
\end{nbwarning}

\subsubsection{Content- and style-based features: LIWC}
\slidesrc{Text features: 10--11}

\textbf{LIWC} (Linguistic Inquiry and Word Count) is a tool with hand-built dictionaries
of psychologically meaningful word categories. It counts, for a text, the percentage of
words in each category: some categories are about style (I-words, function words),
others about content (social words, emotions, cognitive processes). It then combines
them into summary variables.

\begin{nbexample}{LIWC on a technical paragraph}
The slides feed LIWC a paragraph from a computer science paper (about the execution
time of some reasoners analysing log files). The result, compared with the average for
professional or scientific writing:
\begin{center}
\small
\begin{tabular}{@{}l rr@{}}
\toprule
\textbf{LIWC dimension} & \textbf{this text} & \textbf{average (scientific writing)} \\
\midrule
I-words (I, me, my)      & 0.0  & 0.63 \\
social words             & 0.0  & 7.62 \\
positive emotions        & 0.0  & 2.32 \\
negative emotions        & 1.1  & 1.45 \\
cognitive processes      & 17.0 & 7.52 \\
\midrule
\multicolumn{3}{@{}l}{\emph{summary variables}} \\
analytic                 & 99.0 & 92.57 \\
clout                    & 37.5 & 68.17 \\
authenticity             & 65.6 & 24.84 \\
emotional tone           & 11.9 & 43.61 \\
\bottomrule
\end{tabular}
\end{center}
Reading it: the paragraph is extremely \emph{analytic} (99), full of cognitive-process
words (\emph{because}, \emph{depends}, \emph{except}: 17\% vs.\ 7.5\%), impersonal (no
I-words, no social words) and emotionally flat (tone 11.9). This is what we expect from a
technical text, and LIWC finds it by counting words alone.
\end{nbexample}

\subsubsection{Vocabulary and idiosyncrasies}
\slidesrc{Text features: 12}

Features that capture personal or informal writing habits:
\begin{itemize}[nosep]
  \item \textbf{unique words}: the number of distinct words in the sample (vocabulary
        richness);
  \item \textbf{capital words}: the number of words (longer than 2) written only in
        capital letters (SHOUTING);
  \item \textbf{spelling errors}: the number of words not found in the dictionaries of
        the language;
  \item \textbf{character flooding}: \emph{sooooo saaaad\ldots};
  \item \textbf{neologisms}: newly invented words;
  \item \textbf{slang}: \emph{tanta roba!} (Italian slang, ``great stuff!'');
  \item \textbf{emoticons}: \texttt{>(}.
\end{itemize}

\subsubsection{Graph-based features}
\slidesrc{Text features: 13}

A text can also be turned into a \textbf{graph}, and the properties of the graph used as
features. The slide shows an example of \textbf{EmoGraph} for a Spanish sentence
(\emph{He estado tomando cursos en línea sobre temas valiosos que disfruto estudiando y
que podrían ayudarme a hablar en público}, ``I have been taking online courses about
valuable subjects that I enjoy studying and that might help me to speak in public''):
\begin{itemize}[nosep]
  \item the nodes are the \textbf{POS tags} of the words (detailed tags such as
        \texttt{NCMP000}, common noun masculine plural, or \texttt{VMIP1S0}, main verb
        indicative present first person singular) and \textbf{semantic and emotional
        labels} of the words (\emph{education}, \emph{language}, \emph{joy},
        \emph{positive}, \ldots);
  \item the edges follow the order of the words in the text: an edge from a tag to the
        next one, and from a word's tag to its semantic/emotional labels;
  \item graph measures (which nodes are most central, how connected the graph is, \ldots)
        become the features. They capture \emph{how} a person combines grammar and
        emotions, a stylistic fingerprint.
\end{itemize}

\begin{nbsummary}[families of features]
\begin{itemize}[nosep]
  \item Stylistic: n-grams/collocations, POS counts, function words (huge overlap with
        stop words), readability, sentence/word length, punctuation.
  \item Content-based: semantic categories, emotion words, TF-IDF, LSA.
  \item Content + style: LIWC (psychological word categories + summary variables).
  \item Vocabulary and idiosyncrasies: unique words, capitals, spelling errors,
        flooding, neologisms, slang, emoticons.
  \item Graph-based: e.g.\ EmoGraph (POS + emotion labels as a graph).
  \item Scope of a feature: the word (FW), the word in its sentence (POS), the document
        (counts), the corpus (TF-IDF).
\end{itemize}
\end{nbsummary}
